\documentclass[../../main.tex]{subfiles}

\begin{document}

\section{Banach 压缩原理}\label{section:ch10-banach压缩}

\begin{definition}[不动点]\label{definition:ch10-不动点}
$X$ 中的点 $x$ 称为映射 $T:X\to X$ 的不动点, 若 $T(x)=x$.
\end{definition}

这里我们感兴趣的是找到确保映射有不动点的假设. 当然, 映射可以有也可以没有任何不动点. 例如, 映射 $T:\mathbf{R}\to\mathbf{R}$, $T(x)=x+1$ 显然没有不动点.

一元实变函数的不动点对应于使得函数的图像与对角线 $y=x$ 相交的点. 这个观察提供了关于不动点存在性的最为初等的结果的几何解释. 令 $[a,b]$ 为闭有界区间且假定连续函数 $f:[a,b]\to\mathbf{R}$ 的象包含于 $[a,b]$. 则 $f:[a,b]\to\mathbf{R}$ 有一个不动点. 这从介值定理可以得到. 若我们定义 $g(x)=f(x)-x$, 则 $g(a)\geqslant 0$ 且 $g(b)\leqslant 0$, 因此对 $[a,b]$ 中的某个 $x_{0}$, $g(x_{0})=0$, 这意味着 $f(x_{0})=x_{0}$.

$\mathbf{R}^{n}$ 的子集 $K$ 称为是凸的, 若只要 $u$ 和 $v$ 属于 $K$, 线段 $\{tu+(1-t)v\mid 0\leqslant t\leqslant 1\}$ 就包含于 $K$. 前面的结果可推广到欧氏空间 $\mathbf{R}^{n}$ 的子集上的映射: 若 $K$ 是 $\mathbf{R}^{n}$ 的紧凸子集且映射 $T:K\to K$ 是连续的, 则 $T$ 有一个不动点. 这称为 Brouwer 不动点定理. 这里我们将证明称为 Banach 压缩原理的初等的不动点结果. 它对映射做了更多的限制性假设, 但对所处空间的假设非常一般.

\begin{definition}[Lipschitz 映射与压缩]\label{definition:ch10-压缩}
度量空间 $(X,\rho)$ 到其自身的映射 $T$ 称为 Lipschitz 的, 若存在称为映射的 Lipschitz 常数的 $c\geqslant 0$, 使得对所有 $u,v\in X$, $\rho(T(u),T(v))\leqslant c\rho(u,v)$.

若 $c<1$, 该 Lipschitz 映射称为一个压缩.
\end{definition}

\begin{theorem}[Banach 压缩原理]\label{theorem:ch10-banach压缩}
令 $X$ 为完备度量空间而映射 $T:X\to X$ 是一个压缩. 则 $T:X\to X$ 恰有一个不动点.
\end{theorem}

\begin{proof}
令 $c$ 满足 $0\leqslant c<1$, 是映射 $T$ 的 Lipschitz 常数. 选取 $X$ 中的点 $x_{0}$. 现在归纳地定义序列 $\{x_{k}\}$, 首先定义 $x_{1}=T(x_{0})$, 若 $k$ 是一个自然数使得 $x_{k}$ 已被定义, 定义 $x_{k+1}=T(x_{k})$. 序列 $\{x_{n}\}$ 是恰当定义的, 因为 $T(X)$ 是 $X$ 的子集. 我们将证明该序列收敛到 $T$ 的不动点.

事实上, 根据序列的 Lipschitz 常数 $c$ 的定义, 观察到对每个自然数 $k\geqslant 2$,
\begin{align*}
\rho(x_{k+1},x_{k})=\rho(T(x_{k}),T(x_{k-1}))\leqslant c\rho(x_{k},x_{k-1}).
\end{align*}
用归纳法, 得到
\begin{align*}
\rho(x_{k+1},x_{k})\leqslant c^{k}\rho(T(x_{0}),x_{0}).
\end{align*}
因此, 若 $m$ 和 $k$ 是自然数满足 $m>k$, 从关于度量 $\rho$ 的三角不等式和几何求和公式, 得到
\begin{align*}
\rho(x_{m},x_{k})&\leqslant\rho(x_{m},x_{m-1})+\rho(x_{m-1},x_{m-2})+\cdots+\rho(x_{k+1},x_{k})\\
&\leqslant[c^{m-1}+c^{m-2}+\cdots+c^{k}]\rho(T(x_{0}),x_{0})\\
&=c^{k}[1+c+\cdots+c^{m-1-k}]\rho(T(x_{0}),x_{0})\\
&=c^{k}\cdot\frac{1-c^{m-k}}{1-c}\cdot\rho(T(x_{0}),x_{0}).
\end{align*}
因此, 由于 $0\leqslant c<1$,
\begin{align*}
\text{当 }m>k\text{ 时},\quad\rho(x_{m},x_{k})\leqslant\frac{c^{k}}{1-c}\cdot\rho(T(x_{0}),x_{0}).
\end{align*}
但 $\lim_{k\to\infty}c^{k}=0$, 因此从前一个不等式我们得出 $\{x_{k}\}$ 是一个 Cauchy 序列.

根据假设, 度量空间 $X$ 是完备的. 因此存在 $X$ 中的点 $x$ 使得序列 $\{x_{k}\}$ 收敛到它. 由于 $T$ 是 Lipschitz 的, 它是连续的. 因此,
\begin{align*}
T(x)=\lim_{k\to\infty}T(x_{k})=\lim_{k\to\infty}x_{k+1}=x.
\end{align*}
因此映射 $T:X\to X$ 至少有一个不动点. 剩下来要检验它仅有一个不动点. 若 $u$ 和 $v$ 是 $X$ 中的点使得 $T(u)=u$ 与 $T(v)=v$, 则
\begin{align*}
0\leqslant\rho(u,v)=\rho(T(u),T(v))\leqslant c\rho(u,v).
\end{align*}
由于 $0\leqslant c<1$, 我们必须有 $\rho(u,v)=0$, 即 $u=v$. 因此恰有一个不动点.
\end{proof}

以上 Banach 压缩原理的证明实际上不仅仅证明了唯一的不动点的存在性. 它提供了逼近不动点的算法. 事实上, 在 Banach 压缩原理的假设下, 若 $c$ 是映射 $T:X\to X$ 的 Lipschitz 常数, $0\leqslant c<1$, 而 $x_{0}$ 是 $X$ 中的任何点, 则: (i)由 $x_{1}=T(x_{0})$ 和 $k>1$ 时 $x_{k+1}=T(x_{k})$ 递归定义的序列 $\{x_{k}\}$ 收敛到 $T$ 的不动点 $x_{\ast}$; (ii)对每个自然数 $k$, $\rho(x_{\ast},x_{k})\leqslant\frac{c^{k}}{1-c}\cdot\rho(T(x_{0}),x_{0})$.

Banach 压缩原理广泛地用于非线性微分方程的研究中. 我们给出它的应用的一个例子. 假定 $\mathcal{O}$ 是平面 $\mathbf{R}^{2}$ 的包含点 $(x_{0},y_{0})$ 的一个开子集. 给定函数 $g:\mathcal{O}\to\mathbf{R}$, 我们的问题是找到包含 $x_{0}$ 的实数的开区间 $I$ 和可微函数 $f:I\to\mathbf{R}$, 使得
\begin{align}
\text{对所有 }x\in I,\quad f'(x)=g(x,f(x)),\quad f(x_{0})=y_{0}.\label{align:ch10-picard-14}
\end{align}
上述方程的一个非常特殊的情形是 $g$ 与第二个变量无关, 所以 $g(x,y)=h(x)$. 即使在这种情形, 若函数的象 $h:I\to\mathbf{R}$ 不是一个区间, 方程\eqref{align:ch10-picard-14}没有解(见习题 42 和 43). 另一方面, 若 $h$ 是连续的, 则从微积分的基本定理, 由
\begin{align*}
f(x)=y_{0}+\int_{x_{0}}^{x}h(t)\,dt,\quad x\in I
\end{align*}
给出\eqref{align:ch10-picard-14}的唯一解.

因此对一般的二元连续实值函数 $g$, 若连续函数 $f:I\to\mathbf{R}$ 具有性质 $(x,f(x))\in\mathcal{O}$(对每个 $x\in I$), 则 $f$ 是\eqref{align:ch10-picard-14}的解当且仅当对所有 $x\in I$,
\begin{align}
f(x)=y_{0}+\int_{x_{0}}^{x}g(t,f(t))\,dt.\label{align:ch10-picard-15}
\end{align}
如同我们在下一个定理的证明中看到的, 微分方程\eqref{align:ch10-picard-14}与积分方程\eqref{align:ch10-picard-15}的等价性允许我们在微分方程的研究中利用不动点定理.

\begin{theorem}[Picard 局部存在定理]\label{theorem:ch10-picard}
令 $\mathcal{O}$ 为平面 $\mathbf{R}^{2}$ 的包含点 $(x_{0},y_{0})$ 的开子集. 假定函数 $g:\mathcal{O}\to\mathbf{R}^{2}$ 连续且存在正数 $M$ 使得以下 Lipschitz 性质对第二个变量成立:
\begin{align}
\text{对所有 }\mathcal{O}\text{ 中的 }(x,y_{1})\text{ 和 }(x,y_{2}),\quad|g(x,y_{1})-g(x,y_{2})|\leqslant M|y_{1}-y_{2}|.\label{align:ch10-picard-16}
\end{align}
则存在包含 $x_{0}$ 的开区间 $I$ 使得在其上微分方程\eqref{align:ch10-picard-14}有唯一解.
\end{theorem}

\begin{proof}
对正数 $\ell$ 定义 $I_{\ell}$ 为闭区间 $[x_{0}-\ell,x_{0}+\ell]$. 根据上面注意到的\eqref{align:ch10-picard-14}与\eqref{align:ch10-picard-15}的解的等价性, 仅须证明 $\ell$ 可被适当选取, 使得恰好存在一个连续函数 $f:I_{\ell}\to\mathbf{R}$, 它具有性质:
\begin{align*}
\text{对所有 }x\in I_{\ell},\quad f(x)=y_{0}+\int_{x_{0}}^{x}g(t,f(t))\,dt.
\end{align*}
由于 $\mathcal{O}$ 是开的, 我们可以选取正数 $a$ 和 $b$ 使得闭矩形 $R=[x_{0}-a,x_{0}+a]\times[y_{0}-b,y_{0}+b]$ 包含于 $\mathcal{O}$. 现在对每个正数 $\ell$, 定义 $X_{\ell}$ 为度量空间 $C(I_{\ell})$ 的子空间, 具有性质: $|f(x)-y_{0}|\leqslant b$, 对所有 $x\in I_{\ell}$. 即
\begin{align*}
I_{\ell}\text{ 上的连续函数图像包含于矩形 }I_{\ell}\times[y_{0}-b,y_{0}+b]\text{ 中}.
\end{align*}
对于 $f\in X_{\ell}$, 定义函数 $T(f)\in C(I_{\ell})$ 为
\begin{align*}
T(f)(x)=y_{0}+\int_{x_{0}}^{x}g(t,f(t))\,dt,\quad\text{对所有 }x\in I_{\ell}.
\end{align*}
积分方程\eqref{align:ch10-picard-15}的解是映射 $T:X_{\ell}\to C(I_{\ell})$ 的不动点. 证明的策略如下: 由于 $C(I_{\ell})$ 是完备度量空间而 $X_{\ell}$ 是 $C(I_{\ell})$ 的闭子集, $X_{\ell}$ 也是一个完备度量空间. 我们将证明: 若 $\ell$ 取得充分小, 则
\begin{align*}
T(X_{\ell})\subseteq X_{\ell},\quad T:X_{\ell}\to X_{\ell}\text{ 是一个收缩}.
\end{align*}
因此, 我们从 Banach 压缩原理推出 $T:X_{\ell}\to X_{\ell}$ 有唯一的不动点.

为选取 $\ell$ 使得 $T(X_{\ell})\subseteq X_{\ell}$, 我们首先用闭有界矩形 $R$ 的紧性以及 $g$ 的连续性选取一个正数 $K$, 使得对 $R$ 中的所有点 $(x,y)$, $|g(x,y)|\leqslant K$.

现在对 $f\in X_{\ell}$ 和 $x\in I_{\ell}$,
\begin{align*}
|T(f)(x)-y_{0}|=\left|\int_{x_{0}}^{x}g(t,f(t))\,dt\right|\leqslant\ell K.
\end{align*}
因此当 $\ell K\leqslant b$ 时, $T(X_{\ell})\subseteq X_{\ell}$.

观察到对函数 $f_{1},f_{2}\in X_{\ell}$ 和 $x\in I_{\ell}$, 我们从\eqref{align:ch10-picard-16}推出
\begin{align*}
|g(x,f_{1}(x))-g(x,f_{2}(x))|\leqslant M\rho_{\max}(f_{1},f_{2}).
\end{align*}
因此, 用积分的线性和单调性, 我们有
\begin{align*}
|T(f_{1})(x)-T(f_{2})(x)|&=\left|\int_{x_{0}}^{x}[g(t,f_{1}(t))-g(t,f_{2}(t))]\,dt\right|\\
&\leqslant|x-x_{0}|M\rho_{\max}(f_{1},f_{2})\leqslant\ell M\rho_{\max}(f_{1},f_{2}).
\end{align*}
当 $\ell K\leqslant b$ 时, 该不等式以及包含关系 $T(X_{\ell})\subseteq X_{\ell}$ 蕴涵
\begin{align*}
\text{若 }\ell K\leqslant b\text{ 且 }\ell M\leqslant 1,\quad T:X_{\ell}\to X_{\ell}\text{ 是一个压缩}.
\end{align*}
定义 $\ell=\min\{b/K,1/2M\}$. Banach 压缩原理告诉我们映射 $T:X_{\ell}\to X_{\ell}$ 有唯一的不动点.
\end{proof}

\begin{exercise}\label{exercise:ch10-29}
令 $p$ 为多项式. 证明 $p:\mathbf{R}\to\mathbf{R}$ 是 Lipschitz 的当且仅当 $p$ 的次数小于 $2$.
\end{exercise}

\begin{exercise}\label{exercise:ch10-30}
固定 $a>0$. 对 $[0,1]$ 中的 $x$ 定义 $f(x)=ax(1-x)$.
\begin{enumerate}[(1)]
\item $a$ 为何值时 $f([0,1])\subseteq[0,1]$?
\item $a$ 为何值时 $f([0,1])\subseteq[0,1]$ 且 $f:[0,1]\to[0,1]$ 是一个收缩?
\end{enumerate}
\end{exercise}

\begin{exercise}\label{exercise:ch10-31}
度量空间到其自身的 Lipschitz 常数小于 $1$ 的 Lipschitz 映射必有一个不动点吗?
\end{exercise}

\begin{exercise}\label{exercise:ch10-32}
完备度量空间到其自身的 Lipschitz 常数为 $1$ 的 Lipschitz 映射必有一个不动点吗?
\end{exercise}

\begin{exercise}\label{exercise:ch10-33}
令 $X$ 为紧度量空间而 $T$ 是从 $X$ 到其自身的映射, 使得对所有 $u,v\in X$,
\begin{align*}
\rho(T(u),T(v))<\rho(u,v).
\end{align*}
证明 $T$ 有唯一的不动点.
\end{exercise}

\begin{exercise}\label{exercise:ch10-34}
对所有实数 $x$ 定义 $f(x)=\pi/2+x-\arctan x$. 证明:
\begin{align*}
\text{对所有 }u,v\in\mathbf{R},\quad|f(u)-f(v)|<|u-v|.
\end{align*}
证明 $f$ 没有不动点. 这与前一个习题矛盾吗?
\end{exercise}

\begin{exercise}\label{exercise:ch10-35}
在欧氏空间 $\mathbf{R}^{n}$ 中考虑闭单位球 $B=\{x\in\mathbf{R}^{n}\mid\|x\|\leqslant 1\}$. 令 $f$ 为将 $B$ 映到 $B$ 的 Lipschitz 常数为 $1$ 的 Lipschitz 映射. 不用 Brouwer 不动点定理, 证明 $f$ 有一个不动点.
\end{exercise}

\begin{exercise}\label{exercise:ch10-36}
假定映射 $f:\mathbf{R}^{n}\to\mathbf{R}^{n}$ 是一个压缩. 对 $\mathbf{R}^{n}$ 中的所有 $x$, 定义 $g(x)=x-f(x)$. 证明映射 $g:\mathbf{R}^{n}\to\mathbf{R}^{n}$ 既是一对一又是映上的. 也证明 $g$ 和它的逆是连续的.
\end{exercise}

\begin{exercise}\label{exercise:ch10-37}
令 $X$ 为包含点 $x_{0}$ 的完备度量空间且令 $r$ 为正实数. 定义 $K=\{x\in X\mid\rho(x,x_{0})\leqslant r\}$. 假定映射 $T:K\to X$ 是 Lipschitz 常数为 $c$ 的 Lipschitz 映射, 且 $cr+\rho(T(x_{0}),x_{0})\leqslant r$. 证明 $T(K)\subseteq K$ 且 $T:K\to X$ 有一个不动点.
\end{exercise}

\begin{exercise}\label{exercise:ch10-38}
证明: 若函数 $g:\mathbf{R}^{2}\to\mathbf{R}$ 有连续的一阶偏导数, 则对 $\mathbf{R}^{2}$ 中的每个点 $(x_{0},y_{0})$ 存在 $(x_{0},y_{0})$ 的邻域 $\mathcal{O}$ 使得在其上 Lipschitz 假设\eqref{align:ch10-picard-16}成立.
\end{exercise}

\begin{exercise}\label{exercise:ch10-39}
函数 $g:\mathbf{R}^{2}\to\mathbf{R}$ 具有形式 $g(x,y)=h(x)+by$, 其中函数 $h:\mathbf{R}\to\mathbf{R}$ 是连续的, 证明以下是\eqref{align:ch10-picard-14}的解的直接表达式:
\begin{align*}
\text{对 }I\text{ 中的所有 }x,\quad f(x)=e^{b(x-x_{0})}y_{0}+\int_{x_{0}}^{x}e^{b(x-t)}h(t)\,dt.
\end{align*}
\end{exercise}

\begin{exercise}\label{exercise:ch10-40}
考虑微分方程
\begin{align*}
\text{对所有 }x\in\mathbf{R},\quad f'(x)=3[f(x)]^{2/3},\quad f(0)=0.
\end{align*}
证明恒等于零的函数 $f:\mathbf{R}\to\mathbf{R}$ 是解, 且定义为若 $x<0$ 有 $f(x)=0$ 而若 $x\geqslant 0$ 有 $f(x)=x^{3}$ 的函数 $f:\mathbf{R}\to\mathbf{R}$ 也是解. 这与 Picard 存在定理矛盾吗?
\end{exercise}

\begin{exercise}\label{exercise:ch10-41}
对于正数 $\varepsilon$, 考虑微分方程
\begin{align*}
f'(x)=(1/\varepsilon)[1+(f(x))^{2}],\quad\text{对所有 }x\in\mathbf{R},\quad f(0)=0.
\end{align*}
证明在区间 $I=(-\varepsilon(\pi/2),\varepsilon(\pi/2))$ 上存在该微分方程的唯一解 $f(x)=\tan(x/\varepsilon)$ 且在严格包含 $I$ 的区间上不存在解.
\end{exercise}

\begin{exercise}\label{exercise:ch10-42}
令 $I$ 为 $\mathbf{R}$ 中的开区间, 假定函数 $h:I\to\mathbf{R}$ 具有性质: 存在 $I$ 中的点 $x_{1}<x_{2}$ 以及数 $c$ 使得 $h(x_{1})<c<h(x_{2})$ 但 $c$ 不属于 $h(I)$. 通过证明若 $f:I\to\mathbf{R}$ 是一个解, 则连续函数 $f(x)-cx$ 在区间 $[x_{1},x_{2}]$ 上取不到最小值, 证明微分方程\eqref{align:ch10-picard-14}没有解.
\end{exercise}

\begin{exercise}\label{exercise:ch10-43}
用前一个习题叙述与证明 Darboux 的以下定理: 令 $I$ 为 $\mathbf{R}$ 中的开区间, 假定 $f:I\to\mathbf{R}$ 是可微的. 则导数 $f':I\to\mathbf{R}$ 的象是一个区间.
\end{exercise}

\begin{exercise}\label{exercise:ch10-44}
叙述和证明以下微分方程系统的 Picard 存在性定理: $\mathcal{O}$ 是 $\mathbf{R}\times\mathbf{R}^{n}$ 的开子集, $g:\mathcal{O}\to\mathbf{R}^{n}$ 是连续的, 点 $(x_{0},y_{0})$ 在 $\mathcal{O}$ 中, 微分方程系统是:
\begin{align*}
f'(x)=g(x,f(x)),\quad\text{对所有 }x\in I,\quad f(x_{0})=y_{0}.
\end{align*}
(提示: 用 Lipschitz 映射逼近 $g$, 接着用 Arzel\`a-Ascoli 定理.)
\end{exercise}

\end{document}